\scanpage{9}
% Typed face read after a 180-degree rotation.
\uline{Variétés de groupe}\hfill - 22 -\hfill Tribu n° 62

a) à un hom. d'alg. top. augm. de $\hat P(M)$ dans $\hat P(M)\mathbin{\widehat\otimes}\hat P(M)$

b) à un hom. de coalg. de $U(M)\otimes U(M)$ dans $U(M)$ ; i.e. à une mult. bilin. sur $U(M)$, pour laquelle $\Delta:U(M)\longrightarrow U(M)\otimes U(M)$ soit un hom.

On ne sait si l'on définira en général la notion d'hyperalgèbre.

P. 3-13 : notation $d^{(\unsure{\sigma|\gamma})}_{\alpha,\beta}$ pour éviter la confusion avec les monômes. La formule (30) exprime que le gradué associé de $U(G)$ est $\Gamma(T(G))$.

P. 3-14 : on pense qu'il y a trois niveaux pour les calculs : le niveau intrinsèque (par ex. calculs dans les groupes $G_A$), le niveau semi-intrinsèque (on a un module $M$ et $m$ dans $\hat P(M\times M;M)$), enfin le niveau des coordonnées (on raisonne sur des séries formelles). La formule (31) et les ex. 1 et 2 sont au niveau semi-intrinsèque ; si $m(x,y)\equiv x+y+B(x,y)$ mod. deg. 2, l'alg. de Lie est $M$ avec le crochet $[x,y]=B(x,y)-B(y,x)$ ; en particulier dans $B_m$, on a $m(x,y)=x+y+xy$, d'où $[x,y]=xy-yx$. L'alg. de convolution du gr. additif de $M$ est $(M)$.

La déf. d'un morphisme formel au moyen de coordonnées doit être fait au chap. des variétés formelles.

P. 3-15 à 3-19 : Les groupes $G(A)$ viennent trop tard. En alg. IV, on saura substituer des élts. top. nilpotents tendant vers 0 dans une série formelle, et des élts qconques dans une série restreinte. Regrouper les n° 8 et 9 ; supprimer lignes 1 à 4 haut et 1 à 7 bas. Remplacer $a_k$ par $b_k$ dans (36) et (37).

On pourrait donner les sorites suivants pour les noyaux formels :

a) Bijection $\underline V(A)\simeq\operatorname{Hom}_{\mathrm{alg.top.}}(\hat P_+(\underline V),A)$ pour $A$ alg. top. comm. (pas d'unité exigée) ; la top. sur $\hat P_+(\underline V)$ est la top. faible pour la dualité avec $U(\underline V)$.

b) Tout hom. continu $\sigma:A\longrightarrow B$ donne un $\underline V(\sigma):\underline V(A)\longrightarrow\underline V(B)$.

c) Tout $f:\underline V\longrightarrow\underline W$ donne des $f_A:\underline V(A)\longrightarrow\underline W(A)$.

d) Bijection canonique $(\underline V\times\underline W)(A)=\underline V(A)\times\underline W(A)$.

e) On a $\underline V(A)=\lim_J\underline V(A/J)$ ($J$ idéal ouvert) par constr.

f) Lorsque $A^2=0$, bijection $\underline V(A)\simeq T(\underline V)\mathbin{\widehat\otimes} A$.

Si $G$ est un groupe formel, $G(A)$ est un gr. et $G(\ )$ un hom. de groupes.

g) Si $J$ est un idéal fermé de $A$, on a une suite exacte :
\diagram{\begin{tikzcd}
1 \arrow[r] & G_J \arrow[r] & G_A \arrow[r] & G_{A/J} \arrow[r,dashed] & 1
\end{tikzcd}}
(dernière flèche slt si $\dim G$ finie ou $J$ ouvert). Ajouter prop. 4, (i), (ii).
\mspage{8}
